% Options for packages loaded elsewhere
\PassOptionsToPackage{unicode}{hyperref}
\PassOptionsToPackage{hyphens}{url}
\documentclass[
  12pt,
]{article}
\usepackage{xcolor}
\usepackage[a4paper,left=3cm,right=2cm,top=3cm,bottom=2cm]{geometry}
\usepackage{amsmath,amssymb}
\setcounter{secnumdepth}{5}
\usepackage{iftex}
\ifPDFTeX
  \usepackage[T1]{fontenc}
  \usepackage[utf8]{inputenc}
  \usepackage{textcomp} % provide euro and other symbols
\else % if luatex or xetex
  \usepackage{unicode-math} % this also loads fontspec
  \defaultfontfeatures{Scale=MatchLowercase}
  \defaultfontfeatures[\rmfamily]{Ligatures=TeX,Scale=1}
\fi
\usepackage{lmodern}
\ifPDFTeX\else
  % xetex/luatex font selection
  \setmainfont[]{Times New Roman}
\fi
% Use upquote if available, for straight quotes in verbatim environments
\IfFileExists{upquote.sty}{\usepackage{upquote}}{}
\IfFileExists{microtype.sty}{% use microtype if available
  \usepackage[]{microtype}
  \UseMicrotypeSet[protrusion]{basicmath} % disable protrusion for tt fonts
}{}
\usepackage{setspace}
\makeatletter
\@ifundefined{KOMAClassName}{% if non-KOMA class
  \IfFileExists{parskip.sty}{%
    \usepackage{parskip}
  }{% else
    \setlength{\parindent}{0pt}
    \setlength{\parskip}{6pt plus 2pt minus 1pt}}
}{% if KOMA class
  \KOMAoptions{parskip=half}}
\makeatother
\usepackage{longtable,booktabs,array}
\newcounter{none} % for unnumbered tables
\usepackage{calc} % for calculating minipage widths
% Correct order of tables after \paragraph or \subparagraph
\usepackage{etoolbox}
\makeatletter
\patchcmd\longtable{\par}{\if@noskipsec\mbox{}\fi\par}{}{}
\makeatother
% Allow footnotes in longtable head/foot
\IfFileExists{footnotehyper.sty}{\usepackage{footnotehyper}}{\usepackage{footnote}}
\makesavenoteenv{longtable}
\usepackage{graphicx}
\makeatletter
\newsavebox\pandoc@box
\newcommand*\pandocbounded[1]{% scales image to fit in text height/width
  \sbox\pandoc@box{#1}%
  \Gscale@div\@tempa{\textheight}{\dimexpr\ht\pandoc@box+\dp\pandoc@box\relax}%
  \Gscale@div\@tempb{\linewidth}{\wd\pandoc@box}%
  \ifdim\@tempb\p@<\@tempa\p@\let\@tempa\@tempb\fi% select the smaller of both
  \ifdim\@tempa\p@<\p@\scalebox{\@tempa}{\usebox\pandoc@box}%
  \else\usebox{\pandoc@box}%
  \fi%
}
% Set default figure placement to htbp
\def\fps@figure{htbp}
\makeatother
\setlength{\emergencystretch}{3em} % prevent overfull lines
\providecommand{\tightlist}{%
  \setlength{\itemsep}{0pt}\setlength{\parskip}{0pt}}
\usepackage{indentfirst}
\usepackage{booktabs}
\usepackage{longtable}
\usepackage{caption}
\setlength{\parindent}{1.25cm}
\setlength{\parskip}{0pt}
\captionsetup{font=small,labelfont=bf,justification=centering}
\usepackage{bookmark}
\IfFileExists{xurl.sty}{\usepackage{xurl}}{} % add URL line breaks if available
\urlstyle{same}
\hypersetup{
  pdftitle={Dengue e temperatura media diaria em Campos dos Goytacazes{,} RJ},
  hidelinks,
  pdfcreator={LaTeX via pandoc}}

\title{Dengue e temperatura media diaria em Campos dos Goytacazes, RJ}
\author{}
\date{\vspace{-2.5em}}

\begin{document}
\maketitle

\setstretch{1.5}
\section{Materiais}\label{materiais}

Foram utilizados arquivos locais de notificacoes de dengue em formato
Excel (\texttt{DENGUE2020.xlsx} a \texttt{DENGUE2025.xlsx}) e a base
\texttt{temperatura\_diaria\_campos.csv}, contendo temperatura media
diaria para Campos dos Goytacazes, RJ. A unidade geografica analisada
foi o municipio de Campos dos Goytacazes, RJ, codigo IBGE 330100.

A variavel de desfecho foi o numero diario de notificacoes de dengue,
agregado pela data dt\_notific. A variavel de exposicao foi a
temperatura media diaria em graus Celsius. O periodo comum analisado foi
de 2020-01-08 a 2025-12-30. Todos os arquivos foram mantidos localmente.

\section{Metodologia}\label{metodologia}

A analise foi conduzida em R 4.6.0. As bases anuais de dengue foram
padronizadas, os codigos municipais foram normalizados e os registros
foram agregados por dia. A base meteorologica foi lida em formato CSV,
com conversao das variaveis de temperatura para formato numerico.

Foi utilizada exclusivamente correlacao de Spearman. O coeficiente foi
escolhido por ser adequado a contagens assimetricas, com muitos zeros e
valores extremos, e por avaliar associacao monotonica sem pressupor
relacao linear. Foram estimadas associacoes diaria, mensal, anual
exploratoria, por anomalias mensais e por defasagem de temperatura de 0
a 45 dias. O melhor lag foi definido pelo maior valor absoluto de rho de
Spearman.

Os graficos foram exportados apenas como imagens PNG individuais, em
estilo editorial com paleta propria inspirada em Cell Press. Cada
grafico apresenta diretamente o rho de Spearman e o p-valor
correspondente.

\section{Resultados}\label{resultados}

\begin{longtable}[]{@{}
  >{\raggedright\arraybackslash}p{(\linewidth - 12\tabcolsep) * \real{0.1096}}
  >{\raggedright\arraybackslash}p{(\linewidth - 12\tabcolsep) * \real{0.0890}}
  >{\raggedright\arraybackslash}p{(\linewidth - 12\tabcolsep) * \real{0.1507}}
  >{\raggedright\arraybackslash}p{(\linewidth - 12\tabcolsep) * \real{0.1644}}
  >{\raggedright\arraybackslash}p{(\linewidth - 12\tabcolsep) * \real{0.1575}}
  >{\raggedright\arraybackslash}p{(\linewidth - 12\tabcolsep) * \real{0.1644}}
  >{\raggedright\arraybackslash}p{(\linewidth - 12\tabcolsep) * \real{0.1644}}@{}}
\caption{Resumo descritivo da base diaria analisada.}\tabularnewline
\toprule\noalign{}
\begin{minipage}[b]{\linewidth}\raggedright
Dias analisados
\end{minipage} & \begin{minipage}[b]{\linewidth}\raggedright
Casos totais
\end{minipage} & \begin{minipage}[b]{\linewidth}\raggedright
Media diaria de casos
\end{minipage} & \begin{minipage}[b]{\linewidth}\raggedright
Mediana diaria de casos
\end{minipage} & \begin{minipage}[b]{\linewidth}\raggedright
Maximo diario de casos
\end{minipage} & \begin{minipage}[b]{\linewidth}\raggedright
Dias com zero casos (\%)
\end{minipage} & \begin{minipage}[b]{\linewidth}\raggedright
Temperatura media geral
\end{minipage} \\
\midrule\noalign{}
\endfirsthead
\toprule\noalign{}
\begin{minipage}[b]{\linewidth}\raggedright
Dias analisados
\end{minipage} & \begin{minipage}[b]{\linewidth}\raggedright
Casos totais
\end{minipage} & \begin{minipage}[b]{\linewidth}\raggedright
Media diaria de casos
\end{minipage} & \begin{minipage}[b]{\linewidth}\raggedright
Mediana diaria de casos
\end{minipage} & \begin{minipage}[b]{\linewidth}\raggedright
Maximo diario de casos
\end{minipage} & \begin{minipage}[b]{\linewidth}\raggedright
Dias com zero casos (\%)
\end{minipage} & \begin{minipage}[b]{\linewidth}\raggedright
Temperatura media geral
\end{minipage} \\
\midrule\noalign{}
\endhead
\bottomrule\noalign{}
\endlastfoot
2.184 & 38.143 & 17,46 & 1,00 & 458 & 49,2 & 23,84 \\
\end{longtable}

\begin{longtable}[]{@{}
  >{\raggedright\arraybackslash}p{(\linewidth - 10\tabcolsep) * \real{0.4598}}
  >{\raggedright\arraybackslash}p{(\linewidth - 10\tabcolsep) * \real{0.1034}}
  >{\raggedleft\arraybackslash}p{(\linewidth - 10\tabcolsep) * \real{0.0575}}
  >{\raggedright\arraybackslash}p{(\linewidth - 10\tabcolsep) * \real{0.1839}}
  >{\raggedright\arraybackslash}p{(\linewidth - 10\tabcolsep) * \real{0.0920}}
  >{\raggedleft\arraybackslash}p{(\linewidth - 10\tabcolsep) * \real{0.1034}}@{}}
\caption{Correlacoes de Spearman estimadas.}\tabularnewline
\toprule\noalign{}
\begin{minipage}[b]{\linewidth}\raggedright
Analise
\end{minipage} & \begin{minipage}[b]{\linewidth}\raggedright
Metodo
\end{minipage} & \begin{minipage}[b]{\linewidth}\raggedleft
n
\end{minipage} & \begin{minipage}[b]{\linewidth}\raggedright
Rho de Spearman
\end{minipage} & \begin{minipage}[b]{\linewidth}\raggedright
p-valor
\end{minipage} & \begin{minipage}[b]{\linewidth}\raggedleft
lag\_dias
\end{minipage} \\
\midrule\noalign{}
\endfirsthead
\toprule\noalign{}
\begin{minipage}[b]{\linewidth}\raggedright
Analise
\end{minipage} & \begin{minipage}[b]{\linewidth}\raggedright
Metodo
\end{minipage} & \begin{minipage}[b]{\linewidth}\raggedleft
n
\end{minipage} & \begin{minipage}[b]{\linewidth}\raggedright
Rho de Spearman
\end{minipage} & \begin{minipage}[b]{\linewidth}\raggedright
p-valor
\end{minipage} & \begin{minipage}[b]{\linewidth}\raggedleft
lag\_dias
\end{minipage} \\
\midrule\noalign{}
\endhead
\bottomrule\noalign{}
\endlastfoot
Diaria: temperatura media vs casos & spearman & 2184 & 0,118 &
\textless{} 0,001 & NA \\
Anomalias mensais: temperatura vs casos & spearman & 2184 & 0,190 &
\textless{} 0,001 & NA \\
Mensal: temperatura media vs casos & spearman & 72 & 0,155 & 0,195 &
NA \\
Anual: temperatura media vs casos & spearman & 6 & 0,829 & 0,042 & NA \\
Melhor defasagem diaria & spearman & 2142 & 0,189 & \textless{} 0,001 &
42 \\
\end{longtable}

A correlacao diaria bruta foi positiva e fraca (rho = 0,118, p
\textless{} 0,001). A maior correlacao em modulo foi observada com lag
de 42 dias (rho = 0,189, p \textless{} 0,001). A tabela diaria completa
foi salva em \texttt{tabelas/tabela\_diaria\_associacao.csv}; abaixo
esta uma amostra inicial.

\begin{longtable}[]{@{}
  >{\raggedright\arraybackslash}p{(\linewidth - 12\tabcolsep) * \real{0.1594}}
  >{\raggedleft\arraybackslash}p{(\linewidth - 12\tabcolsep) * \real{0.0870}}
  >{\raggedright\arraybackslash}p{(\linewidth - 12\tabcolsep) * \real{0.1739}}
  >{\raggedleft\arraybackslash}p{(\linewidth - 12\tabcolsep) * \real{0.1304}}
  >{\raggedright\arraybackslash}p{(\linewidth - 12\tabcolsep) * \real{0.1449}}
  >{\raggedright\arraybackslash}p{(\linewidth - 12\tabcolsep) * \real{0.1884}}
  >{\raggedright\arraybackslash}p{(\linewidth - 12\tabcolsep) * \real{0.1159}}@{}}
\caption{Amostra da tabela diaria com casos, temperatura, rho de
Spearman e p-valor.}\tabularnewline
\toprule\noalign{}
\begin{minipage}[b]{\linewidth}\raggedright
Data
\end{minipage} & \begin{minipage}[b]{\linewidth}\raggedleft
Casos
\end{minipage} & \begin{minipage}[b]{\linewidth}\raggedright
Temp. media
\end{minipage} & \begin{minipage}[b]{\linewidth}\raggedleft
Lag dias
\end{minipage} & \begin{minipage}[b]{\linewidth}\raggedright
Temp. lag
\end{minipage} & \begin{minipage}[b]{\linewidth}\raggedright
Rho Spearman
\end{minipage} & \begin{minipage}[b]{\linewidth}\raggedright
p-valor
\end{minipage} \\
\midrule\noalign{}
\endfirsthead
\toprule\noalign{}
\begin{minipage}[b]{\linewidth}\raggedright
Data
\end{minipage} & \begin{minipage}[b]{\linewidth}\raggedleft
Casos
\end{minipage} & \begin{minipage}[b]{\linewidth}\raggedright
Temp. media
\end{minipage} & \begin{minipage}[b]{\linewidth}\raggedleft
Lag dias
\end{minipage} & \begin{minipage}[b]{\linewidth}\raggedright
Temp. lag
\end{minipage} & \begin{minipage}[b]{\linewidth}\raggedright
Rho Spearman
\end{minipage} & \begin{minipage}[b]{\linewidth}\raggedright
p-valor
\end{minipage} \\
\midrule\noalign{}
\endhead
\bottomrule\noalign{}
\endlastfoot
2020-01-08 & 1 & 27,95 & 42 & NA & 0,189 & \textless{} 0,001 \\
2020-01-09 & 0 & 27,48 & 42 & NA & 0,189 & \textless{} 0,001 \\
2020-01-10 & 2 & 27,88 & 42 & NA & 0,189 & \textless{} 0,001 \\
2020-01-11 & 0 & 28,35 & 42 & NA & 0,189 & \textless{} 0,001 \\
2020-01-12 & 0 & 28,44 & 42 & NA & 0,189 & \textless{} 0,001 \\
2020-01-13 & 0 & 27,17 & 42 & NA & 0,189 & \textless{} 0,001 \\
2020-01-14 & 0 & 25,41 & 42 & NA & 0,189 & \textless{} 0,001 \\
2020-01-15 & 0 & 25,20 & 42 & NA & 0,189 & \textless{} 0,001 \\
2020-01-16 & 0 & 27,64 & 42 & NA & 0,189 & \textless{} 0,001 \\
2020-01-17 & 0 & 26,54 & 42 & NA & 0,189 & \textless{} 0,001 \\
2020-01-18 & 0 & 24,06 & 42 & NA & 0,189 & \textless{} 0,001 \\
2020-01-19 & 0 & 24,74 & 42 & NA & 0,189 & \textless{} 0,001 \\
2020-01-20 & 0 & 22,68 & 42 & NA & 0,189 & \textless{} 0,001 \\
2020-01-21 & 0 & 23,60 & 42 & NA & 0,189 & \textless{} 0,001 \\
2020-01-22 & 0 & 24,83 & 42 & NA & 0,189 & \textless{} 0,001 \\
2020-01-23 & 0 & 24,48 & 42 & NA & 0,189 & \textless{} 0,001 \\
2020-01-24 & 0 & 22,40 & 42 & NA & 0,189 & \textless{} 0,001 \\
2020-01-25 & 0 & 23,12 & 42 & NA & 0,189 & \textless{} 0,001 \\
2020-01-26 & 0 & 24,40 & 42 & NA & 0,189 & \textless{} 0,001 \\
2020-01-27 & 0 & 24,98 & 42 & NA & 0,189 & \textless{} 0,001 \\
2020-01-28 & 1 & 27,65 & 42 & NA & 0,189 & \textless{} 0,001 \\
2020-01-29 & 0 & 28,58 & 42 & NA & 0,189 & \textless{} 0,001 \\
2020-01-30 & 0 & 28,15 & 42 & NA & 0,189 & \textless{} 0,001 \\
2020-01-31 & 0 & 28,23 & 42 & NA & 0,189 & \textless{} 0,001 \\
2020-02-01 & 0 & 28,68 & 42 & NA & 0,189 & \textless{} 0,001 \\
2020-02-02 & 0 & 27,86 & 42 & NA & 0,189 & \textless{} 0,001 \\
2020-02-03 & 0 & 27,46 & 42 & NA & 0,189 & \textless{} 0,001 \\
2020-02-04 & 0 & 27,16 & 42 & NA & 0,189 & \textless{} 0,001 \\
2020-02-05 & 0 & 27,18 & 42 & NA & 0,189 & \textless{} 0,001 \\
2020-02-06 & 0 & 26,50 & 42 & NA & 0,189 & \textless{} 0,001 \\
\end{longtable}

\begin{figure}

{\centering \includegraphics[width=1\linewidth]{../data_inmet_campos/resultados_dengue_temperatura/figuras/01_serie_diaria_casos_temperatura} 

}

\caption{Serie diaria de notificacoes de dengue e temperatura media diaria.}\label{fig:fig-serie}
\end{figure}

\begin{figure}

{\centering \includegraphics[width=1\linewidth]{../data_inmet_campos/resultados_dengue_temperatura/figuras/02_dispersao_diaria_spearman} 

}

\caption{Associacao diaria bruta entre temperatura media e notificacoes.}\label{fig:fig-dispersao}
\end{figure}

\begin{figure}

{\centering \includegraphics[width=1\linewidth]{../data_inmet_campos/resultados_dengue_temperatura/figuras/03_spearman_por_lag} 

}

\caption{Correlacao de Spearman por defasagem da temperatura media diaria.}\label{fig:fig-lag}
\end{figure}

\begin{figure}

{\centering \includegraphics[width=1\linewidth]{../data_inmet_campos/resultados_dengue_temperatura/figuras/04_agregado_mensal_spearman} 

}

\caption{Associacao mensal entre temperatura media e notificacoes agregadas.}\label{fig:fig-mensal}
\end{figure}

\begin{figure}

{\centering \includegraphics[width=1\linewidth]{../data_inmet_campos/resultados_dengue_temperatura/figuras/05_agregado_anual_spearman} 

}

\caption{Associacao anual exploratoria entre temperatura media anual e notificacoes anuais.}\label{fig:fig-anual}
\end{figure}

\begin{figure}

{\centering \includegraphics[width=1\linewidth]{../data_inmet_campos/resultados_dengue_temperatura/figuras/06_anomalias_mensais_spearman} 

}

\caption{Associacao entre anomalias mensais de temperatura e de notificacoes.}\label{fig:fig-anomalias}
\end{figure}

\end{document}
